\documentclass[11pt,a4paper]{report}
\usepackage{marvosym}

\assignment{1}
\group{73}
\students{Nicolas Heymans}{Gilles de Jamblinne de Meux }

\begin{document}

\maketitle

Answer to the questions by adding text in the boxes. 
You may answer in either \textbf{French or English}. 
Do not modify anything else in the template.  
The size of the boxes indicate the place you \textbf{can} use, but you \textbf{need not} use all of it (it is not an indication of any expected length of answer). 
\textbf{Be as concise as possible! A good short answer is better than a lot of nonsense!}
%\bigskip

\section{Understanding a CP model (2 points)}

\begin{enumerate}
    \item Answer correctly to the questions 1-4. \textbf{(0.5pt)}
    You can help yourself by looking
    at \url{https://www.pycsp.org/documentation/constraints/} to find the definition of the different constraints.
    \item Answer correctly to the question 5. \textbf{(0.5pt)}
    \item Add the needed constraint to the model and submit it on INGInious. \textbf{(1pt)}
\end{enumerate}

\newpage
\section{Minesweeper Model (8 points)}

\begin{enumerate}
    \item Implement your CP model in the file \texttt{minesweeper.py} and submit it on INGInious. (\textbf{4 pts})
    \item Describe your CP model: First describe the variables of your CP model. 
    For each set of variables, explain what they represent, give their domains and how many variables are in this set. 
    Then, list the constraints of your CP model. 
    For each constraint explain what it enforces, and on which variables it is applied. (\textbf{4 pts})
\end{enumerate}

\begin{answers}[10cm]
    \textbf{Variables et Domaines}
    \begin{itemize}
        \item Nous utilisons un tableau de variables \texttt{VarArray} nommé $x$, de taille $rows \times columns$.
        \item Chaque variable $x[i][j]$ a un domaine binaire $\{0, 1\}$: la valeur $1$ indique la présence d'une mine et $0$ son absence.
        \item Le nombre total de variables correspond au nombre total de cellules de la grille, soit $|V| = rows \times columns$.
    \end{itemize}

    \textbf{Contraintes}
    \begin{itemize}
        \item \textbf{Cases d'indices  :} Pour chaque cellule $(i, j)$ contenant un indice $\geq 0$ dans la grille d'indice, nous appliquons la contrainte \texttt{satisfy(x[i][j] == 0)} dans la grille de solution, car une mine ne peut pas être placée sur une case d'indice.
        \item \textbf{Somme du voisinage :} Pour chaque indice $C$ situé à la position $(i, j)$, nous imposons que la somme des variables adjacentes $x[ki][kj]$ dans le voisinage de Moore (les 8 cases entourant la cellule, sans s'inclure elle-même) dans la grille de la solution soit exactement égale à $C$.   
        \item \textbf{Résolution :} Si le solveur trouve une solution satisfaisable , nous extrayons toutes les coordonnées $(i, j)$ où $x[i][j] = 1$ (les coordonnées des bombes). Dans le cas contraire, le modèle retourne simplement \texttt{None}.
    \end{itemize}
\end{answers}




\newpage
\section{Atom Placement Problem()}
\begin{enumerate}
\item Formulate the atom placement problem as a Local Search problem (problem, cost function, feasible solutions, optimal solutions). \textbf{(1 pt)}
\end{enumerate}

\begin{answers}[6cm]

Trouver la configuration spatiale de $n$ atomes qui minimise l'énergie potentielle totale du système.

Soit un vecteur 
\[
\mathbf{s} = [s_1, s_2, \dots, s_n]
\] 
où $s_i$ est le type d'un atome dans le cristal et $s_i \in \{0, 1, \dots, n-1\}$.

nous cherchons a minimiser $E_s$ qui ce définit comme suit:
\[
E_s = \sum_{(i,j)\in F} H(s_i, s_j)
\]

où $H$ est la matrice $n$ x $n$ contenant les valeurs d'énergie pour une liaison atomique $(s_i, s_j)$.

\end{answers}


\begin{enumerate}
\setcounter{enumi}{1}
    \item You are given a template on Moodle: \textit{atomplacement.py}. Implement your own extension of the \textit{Problem} class from \texttt{aima-python3}. Implement the \texttt{maxvalue} and \texttt{randomized maxvalue} strategies. To do so,
	you can get inspiration from the \textit{randomwalk} function in \textit{search.py}. Your program will be evaluated on 15 instances (during 1 minute) of which 5 are hidden. We expect you to solve at least 12 out the 15.
    \begin{enumerate}
        \item \texttt{maxvalue} chooses the best node (i.e., the node with maximal
        value) in the neighborhood, even if it degrades the quality of the current solution. 
        The \texttt{maxvalue} strategy should be defined in a function called
        \textit{maxvalue} with the following signature: \\
        \texttt{maxvalue(problem,limit=100)}. \textbf{(2.5 pts)}
    \end{enumerate}
\end{enumerate}

\begin{answers}[2.5 cm]
% IN CASE YOU WANT TO SPECIFY ANY COMMENTS ABOUT YOUR IMPLEMENTATION
\end{answers}



\begin{enumerate}
\setcounter{enumi}{1}
\begin{enumerate}
\setcounter{enumii}{1}
    \item \texttt{randomized maxvalue} chooses the next node randomly among
        the 5~best neighbors (again, even if it degrades the quality of the current solution).
        The \texttt{randomized maxvalue} strategy should be defined in a function called
        \textit{randomized\_maxvalue} with the following signature: \\
        \texttt{randomized\_maxvalue(problem,limit=100)}. \textbf{(2.5 pts)}
\end{enumerate}
\end{enumerate}

\begin{answers}[2.5 cm]
% IN CASE YOU WANT TO SPECIFY ANY COMMENTS ABOUT YOUR IMPLEMENTATION
\end{answers}



\begin{enumerate}
\setcounter{enumi}{2}
\item Compare the 2 strategies implemented in the previous question and \texttt{randomwalk} defined in \textit{search.py} on the given vertex cover instances. 
    Run the strategies during 100 steps, and report, in a table, the computation
    time, the value of the best solution and the number of steps needed to reach the best result. For the
    randomized max value and the random walk, each instance should be tested 10~times to reduce
    the effects of the randomness on the result. When multiple runs of the same instance
    are executed, report the mean of the quantities. \textbf{(3 pts)}
\end{enumerate}

\begin{answers}[7cm]
% Your answer here
\begin{center}
\begin{tabular}{||l||l|l|l||l|l|l||l|l|l||}
\hline
\multirow{3}{*}{Inst.} & \multicolumn{3}{c||}{Max value} & \multicolumn{3}{c||}{Random max value} & \multicolumn{3}{c||}{Random walk} \\
\cline{2-10}
\cline{2-10}
 & T(s) & Val & NS & T(s) & Val & NS & T(s) & Val & NS\\
\hline
i01 & 0.2887 & 151.0 & 5.0 & 0.2992 & 129.4 & 59.6 & 0.0160 & 216.3 & 62.1 \\
\hline
i02 & 11.4055 & 1281.0 & 17.0 & 11.4880 & 1251.5 & 44.6 & 0.1115 & 1696.6 & 58.4 \\
\hline
i03 & 13.1580 & 1128.0 & 22.0 & 13.3080 & 1110.3 & 71.4 & 0.1249 & 1834.4 & 64.1 \\
\hline
i04 & 7.0758 & 274.0 & 12.0 & 7.1667 & 259.3 & 47.5 & 0.0721 & 574.5 & 53.3 \\
\hline
i05 & 0.2632 & 304.0 & 4.0 & 0.2717 & 298.3 & 46.5 & 0.0139 & 364.7 & 50.8 \\
\hline
i06 & 0.3199 & 269.0 & 7.0 & 0.3276 & 270.4 & 51.7 & 0.0154 & 339.2 & 47.2 \\
\hline
i07 & 1.6651 & 244.0 & 14.0 & 1.7188 & 224.5 & 61.3 & 0.0361 & 419.7 & 35.1 \\
\hline
i08 & 5.6791 & 465.0 & 20.0 & 5.7849 & 461.1 & 75.7 & 0.0691 & 823.9 & 42.0 \\
\hline
i09 & 1.7871 & 481.0 & 7.0 & 1.7515 & 411.4 & 59.4 & 0.0375 & 688.1 & 27.9 \\
\hline
i10 & 4.8352 & 1994.0 & 16.0 & 4.8834 & 1994.2 & 58.3 & 0.0648 & 2338.8 & 61.9 \\
\hline
\end{tabular}
\end{center}
\textbf{NS}: Number of steps — \textbf{T}: Time — \textbf{Val}: Value of the best solution
\end{answers}



\begin{enumerate}
\setcounter{enumi}{3}
    \item \textbf{(4 pts)} Answer the following questions:
    \begin{enumerate}
        \item In one sentence, what is the best strategy? \textbf{(0.5 pt)}
    \end{enumerate}
\end{enumerate}

\begin{answers}[3.5cm]
La meilleure stratégie est \texttt{Random max value}, car elle a systématiquement des résultats meilleurs ou équivalents à \texttt{Max value}.
\end{answers}


\newpage
\begin{enumerate}
\setcounter{enumi}{3}
\begin{enumerate}
\setcounter{enumii}{1}
    \item Why do you think the best strategy beats the other ones? What conclusions can be drawn on the atom placement problem ? \textbf{(1.5 pt)}
\end{enumerate}
\end{enumerate}

\begin{answers}[7.5cm]
La supériorité de \texttt{Random max value} peut s’expliquer par sa capacité à explorer un plus grand nombre de trajectoires différentes, en sélectionnant aléatoirement une solution parmi un ensemble de candidats. Cette approche lui permet d’éviter de tomber immédiatement dans un minimum local et augmente ses chances d’atteindre un minimum global.

En revanche, \texttt{Max value} ne considère qu’une seule solution à la fois, ce qui limite sa diversification. Ainsi, lorsqu’elle se retrouve dans un minimum local si aucune solution voisine n’améliore la valeur de l’objectif, l’algorithme reste bloqué dans ce minimum local.

À l’opposé, \texttt{Random walk} ne vise pas une solution spécifique : elle ne réalise que de la diversification, sans aucune intensification. Elle se contente d’explorer l’espace de solutions de manière aléatoire, sans objectif précis.

En conclusion, pour le problème de placement d’atomes, les stratégies combinant exploration aléatoire et sélection de solutions jugées prometteuses permettent d’obtenir des solutions de meilleure qualité, proches du minimum global. Elles offrent un équilibre entre diversification et intensification, réduisant le risque de stagnation dans un minimum local et favorisant l’atteinte du minimum global.


\end{answers}



\begin{enumerate}
\setcounter{enumi}{3}
\begin{enumerate}
\setcounter{enumii}{2}
    \item What are the limitations of each strategy in terms of diversification 
    and intensification? \textbf{(1 pt)}
\end{enumerate}
\end{enumerate}

\begin{answers}[5cm]
\textbf{Random Walk}: privilégie uniquement la diversification, sans aucune intensification.\\

\textbf{Max Value}: privilégie uniquement l'intensification, avec presque aucune diversification. \\

\textbf{Random Max Value}:  équilibre entre diversification et intensification (sélection aléatoire parmi les 5 meilleures solutions).\\

\end{answers}

\begin{enumerate}
\setcounter{enumi}{3}
\begin{enumerate}
\setcounter{enumii}{3}
    \item What is the behavior of the different techniques when they fall 
    in a local optimum? \textbf{(1 pt)}
\end{enumerate}
\end{enumerate}
\begin{answers}[5cm]
\textbf{Random Walk}: Lorsqu'il atteint un minimum local trouvé par hasard, l'algorithme le quitte immédiatement en choisissant une nouvelle solution aléatoire.\\


\textbf{Max Value}: Une fois bloqué dans un minimum local, il remonte vers une solution légèrement sous-optimale, puis retourne directement vers le minimum local précédemment trouvé.\\


\textbf{Random Max Value}: Il peut quitter le minimum local et a une certaine probabilité de continuer à remonter afin d’explorer un autre minimum local.\\


\end{answers}


\end{document}
